\section{Benchmarks}

\subsection{QuickStorm: AudioPlayer}

The original QuickStorm paper~\cite{quickstorm} presents a case study on ``TodoMVC,'' a to-do list application. While this example is comprehensive, its complexity makes it less suitable for concise illustration. Instead, we focus on a simpler application—the ``Audio Player''—which is featured in the QuickStorm tutorial.
The Audio Player application allows the user to play or pause audio files by clicking a single button. Each user interaction toggles the state between playing and paused. After each operation, the application displays the current state (either playing or paused) and the elapsed playback time.

\paragraph{PAT model} I have abstracted the application into the following APIs:
\begin{minted}[xleftmargin=5pt, numbersep=4pt, linenos = true, fontsize = \footnotesize, escapeinside=??]{OCaml}
    type init = unit (* initialize the AudioPlayer *)
    type click = unit (* click the button *)
    type displayTime = int (* the elapsed playback time *)
    type displayText = string (* playing or paised *)
\end{minted}
The user can do $\eff{init}$ to initialize the AudioPlayer, and then do $\eff{click}$ to play or pause the audio file. The application will return the current state via $\eff{displayText}$ and the elapsed playback time via $\eff{playTime}$. The property of the AudioPlayer is that the elapsed playback time should be non-decreasing, these we define the following property to test against:
\begin{align*}
    \phi \doteq \exists t_1{:}\Int . \finalA (\msg{displayTime}{\vnu = t_1} \land \nextA \finalA \msg{displayTime}{\vnu < t_2})
\end{align*}\noindent
We further expect executions to adhere to the following behavioral patterns:
\begin{itemize}
    \item \textbf{Paused at initialization}: Upon initialization, the application should be in the paused state.
    \item \textbf{Elapsed time starts at 0}: The displayed playback time should start from 0.
    \item \textbf{Eventual pause}: If the application is playing, it should eventually return to the paused state.
    \item \textbf{Pause preserves display time}: While paused, the displayed playback time should remain unchanged.
    \item \textbf{Automatic time update while playing}: When the application is playing, the displayed playback time should increase automatically.
\end{itemize}

Then, one expected trace is:
\begin{align*}
    &\eff{init}();\eff{displayText}("paused");\\
    &\eff{click}();\eff{displayText}("playing");\eff{displayTime}(0);\eff{displayTime}(1);\eff{displayTime}(0);\\
    &\eff{click}();\eff{displayText}("paused");
\end{align*}\noindent
If click the button too fast, the displayed may not even update.

\begin{figure}[h!]
    \centering
    {\small\begin{align*}
        &\eff{init}: \pat{\emptyset}{\msg{init}{\top}}{\untilA (\neg\msg{click}{\top}) \msg{displayText}{\vnu = \Code{paused}}}
        \\ \\
        &\eff{click}: \pat{\I{LastPaused}}{\msg{click}{\top}}{\msg{displayText}{\vnu = \Code{playing}} \land \untilA (\neg\msg{click}{\top}) \msg{displayTime}{\top}} \interty\\
        &\qquad \pat{\I{LastPaused}}{\msg{click}{\top}}{\msg{displayText}{\vnu = \Code{playing}} \land \untilA (\neg\msg{displayTime}{\top}) \msg{click}{\top}} \interty\\
        &\qquad \pat{\neg\I{LastPaused}}{\msg{click}{\top}}{\msg{displayText}{\vnu = \Code{paused}}}\\
        \\ \\
        &\eff{displayTime}: \pat{\neg\finalA\msg{displayTime}{\top}}{\msg{displayTime}{\vnu = 0}}{\untilA (\neg\msg{click}{\top}) \msg{displayTime}{\vnu \geq 0}} \\
        &\qquad \pat{\neg\finalA\msg{displayTime}{\top}}{\msg{displayTime}{\vnu = 0}}{\untilA (\neg\msg{displayTime}{\top}) \msg{click}{\geq 0}} \interty\\
        &\qquad \pat{\finalA\msg{displayTime}{\top}}{\msg{displayTime}{\vnu \geq 0}}{\untilA (\neg\msg{click}{\top}) \msg{displayTime}{\vnu \geq 0}} \interty\\
        &\qquad \pat{\finalA\msg{displayTime}{\top}}{\msg{displayTime}{\vnu \geq 0}}{\untilA (\neg\msg{displayTime}{\top}) \msg{click}{\top}}
        \\ \\
        &\eff{displayText}: \pat{\anyA^*}{\msg{displayText}{\top}}{\anyA^*}
    \end{align*}}
    \caption{PATs for AudioPlayer}
    \label{fig:audio-player-pat}
\end{figure}

The semantics outlined above are precisely captured by the PATs in Figure~\ref{fig:audio-player-pat}. The $\eff{displayTime}$ event is designed to repeatedly trigger subsequent $\eff{displayTime}$ events, but it can also be interrupted by a $\eff{click}$ event. There are four distinct cases for $\eff{displayTime}$, arising from the combination of two binary choices:
\begin{itemize}
    \item \textbf{Is this the first display time?} The value returned is either $0$ (if it is the first occurrence) or a natural number, as determined by the history automaton $\finalA\msg{displayTime}{\top}$.
    \item \textbf{Is the event interpreted by a click?} The prophecy automaton specifies either that ``no click occurs before the next display time'' ($\untilA (\neg\msg{click}{\top}) \msg{displayTime}{\vnu \geq 0}$), or that ``a click must occur before any futher display time happens'' ($\untilA (\neg\msg{displayTime}{\top}) \msg{click}{\top}$).
\end{itemize}

\begin{figure}[h!]
    \centering
    \begin{minted}[xleftmargin=5pt, numbersep=4pt, linenos = true, fontsize = \footnotesize, escapeinside=??]{OCaml}
    ?$\mykw{gen}$? ?$\eff{init}$? in
    let (txt1: string) = ?$\mykw{obs}$? ?$\eff{displayText}$? in ?$\mykw{assert}$? (txt1 == "paused") in 
    ?$\mykw{gen}$? ?$\eff{click}$? in
    let (txt2: string) = ?$\mykw{obs}$? ?$\eff{displayText}$? in ?$\mykw{assert}$? (txt2 == "playing") in 
    let (t1: int) = ?$\mykw{obs}$? ?$\eff{displayTime}$? in ?$\mykw{assert}$? (t1 == 0) in
    loop (let (t2: int) = ?$\mykw{obs}$? ?$\eff{displayTime}$? in ?$\mykw{assert}$? (t2 >= 0)); 
    let (t3: int) = ?$\mykw{obs}$? ?$\eff{displayTime}$? in ?$\mykw{assert}$? (t3 >= 0) in 
    loop (let (t4: int) = ?$\mykw{obs}$? ?$\eff{displayTime}$? in ?$\mykw{assert}$? (t4 >= 0));
    let (t5: int) = ?$\mykw{obs}$? ?$\eff{displayTime}$? in ?$\mykw{assert}$? (t5 < t3) in
    ?$\mykw{gen}$? ?$\eff{click}$? in
    let (txt3: string) = ?$\mykw{obs}$? ?$\eff{displayText}$? in ?$\mykw{assert}$? (txt3 == "paused")
    \end{minted}
    \caption{A test generator for audio player }\label{fig:audio-player-gen}
\end{figure}
    
Figure~\ref{fig:audio-player-gen} shows a synthesized test generator for the audio player. The generator begins by initializing the application with $\eff{init}$ (line 1). It then simulates playing the audio and observes the displayed time values, ensuring that the sequence of displayed times is non-decreasing up to a certain point: specifically, it checks that there exists a displayed time $t_3$ that is greater than a subsequently displayed time $t_5$, with both $t_3$ and $t_5$ potentially preceded by arbitrarily many displayed times (captured by the two loops). Finally, the generator simulates clicking the button to pause the audio.